% Bibliography of the revised manuscript. Every entry verified (paper/REFERENCES_AUDIT.md for the 29 retained
% original entries, paper/REFERENCES_VERIFIED.md for the new ones). Two submitted entries were removed as unverifiable.
\begin{thebibliography}{99}

\bibitem{preskill2018quantum}
Preskill, J.: Quantum computing in the NISQ era and beyond. Quantum \textbf{2}, 79 (2018). \doi{10.22331/q-2018-08-06-79}

\bibitem{arute2019quantum}
Arute, F., Arya, K., Babbush, R., Bacon, D., Bardin, J.C., Barends, R., \emph{et al.}: Quantum supremacy using a programmable superconducting processor. Nature \textbf{574}(7779), 505--510 (2019). \doi{10.1038/s41586-019-1666-5}

\bibitem{pednault2017breaking}
Pednault, E., Gunnels, J.A., Nannicini, G., Horesh, L., Magerlein, T., Solomonik, E., Draeger, E.W., Holland, E.T., Wisnieff, R.: Pareto-efficient quantum circuit simulation using tensor contraction deferral (first version titled ``Breaking the 49-qubit barrier in the simulation of quantum circuits''). arXiv:1710.05867 (2017)

\bibitem{villalonga2019flexible}
Villalonga, B., Boixo, S., Nelson, B., Henze, C., Rieffel, E., Biswas, R., Mandr\`a, S.: A flexible high-performance simulator for verifying and benchmarking quantum circuits implemented on real hardware. npj Quantum Information \textbf{5}, 86 (2019). \doi{10.1038/s41534-019-0196-1}

\bibitem{jones2019quest}
Jones, T., Brown, A., Bush, I., Benjamin, S.C.: QuEST and high performance simulation of quantum computers. Scientific Reports \textbf{9}, 10736 (2019). \doi{10.1038/s41598-019-47174-9}

\bibitem{suzuki2021qulacs}
Suzuki, Y., \emph{et al.}: Qulacs: a fast and versatile quantum circuit simulator for research purpose. Quantum \textbf{5}, 559 (2021). \doi{10.22331/q-2021-10-06-559}

\bibitem{khalid2004fpga}
Khalid, A.U., Zilic, Z., Radecka, K.: FPGA emulation of quantum circuits. In: Proc. IEEE International Conference on Computer Design (ICCD), pp. 310--315 (2004). \doi{10.1109/ICCD.2004.1347938}

\bibitem{lee2016fpga}
Lee, Y.H., Khalil-Hani, M., Marsono, M.N.: An FPGA-based quantum computing emulation framework based on serial-parallel architecture. International Journal of Reconfigurable Computing \textbf{2016}, 5718124 (2016). \doi{10.1155/2016/5718124}

\bibitem{aminian2008fpga}
Aminian, M., Saeedi, M., Zamani, M.S., Sedighi, M.: FPGA-based circuit model emulation of quantum algorithms. In: Proc. IEEE Computer Society Annual Symposium on VLSI (ISVLSI) (2008). \doi{10.1109/ISVLSI.2008.43}

\bibitem{mahmud2019scaling}
Mahmud, N., El-Araby, E., Caliga, D.: Scaling reconfigurable emulation of quantum algorithms at high precision and high throughput. Quantum Engineering \textbf{1}(2), e19 (2019). \doi{10.1002/que2.19}

\bibitem{pilch2019fpga}
Pilch, J., D{\l}ugopolski, J.: An FPGA-based real quantum computer emulator. Journal of Computational Electronics \textbf{18}, 329--342 (2019). \doi{10.1007/s10825-018-1287-5}

\bibitem{moawad2022cfd}
Moawad, Y., Vanderbauwhede, W., Steijl, R.: Investigating hardware acceleration for simulation of CFD quantum circuits. Frontiers in Mechanical Engineering \textbf{8}, 925637 (2022). \doi{10.3389/fmech.2022.925637}

\bibitem{moawad2024iteration}
Moawad, Y., Brown, A., Steijl, R., Vanderbauwhede, W.: Optimising iteration scheduling for full-state vector simulation of quantum circuits on FPGAs. arXiv:2411.18354 (2024)

\bibitem{suzuki2023hybrid}
Suzuki, T., Miyazaki, T., Inaritai, T., Otsuka, T.: Quantum AI simulator using a hybrid CPU--FPGA approach. Scientific Reports \textbf{13}, 7735 (2023). \doi{10.1038/s41598-023-34600-2}

\bibitem{bennakhi2025q2sv}
Bennakhi, A., Byrd, G.T., Franzon, P.: Q2SV: A high-level synthesis approach for state vector quantum simulation. Quantum Engineering \textbf{2025}, 9017796 (2025). \doi{10.1155/que2/9017796}

\bibitem{tran2025hpqea}
Tran, V.D., Vu, T.H., Le, V.T.D., Pham, H.L., Nakashima, Y.: HPQEA: A scalable and high-performance quantum emulator with high-bandwidth memory for diverse algorithms support. arXiv:2510.07110 (2025)

\bibitem{cicero2024review}
Cicero, A., Maleki, M.A., Azhar, M.W., Kockum, A.F., Trancoso, P.: Simulation of quantum computers: review and acceleration opportunities. arXiv:2410.12660 (2024)

\bibitem{tran2025qea}
Tran, V.D., Vu, T.H., Le, V.T.D., Pham, H.L., Nakashima, Y.: QEA: An accelerator for quantum circuit simulation with resources efficiency and flexibility. In: Proc. 10th International Conference on ICs, Design and Verification (ICDV) (2025). arXiv:2503.14951

\bibitem{wang2020shuhai}
Wang, Z., Huang, H., Zhang, J., Alonso, G.: Benchmarking high bandwidth memory on FPGAs. arXiv:2005.04324 (2020)

\bibitem{choi2020hlshbm}
Choi, Y.-k., Chi, Y., Wang, J., Guo, L., Cong, J.: When HLS meets FPGA HBM: Benchmarking and bandwidth optimization. arXiv:2010.06075 (2020)

\bibitem{choi2021hbmconnect}
Choi, Y.-k., Chi, Y., Qiao, W., Samardzic, N., Cong, J.: HBM Connect: High-performance HLS interconnect for FPGA HBM. In: Proc. ACM/SIGDA International Symposium on Field-Programmable Gate Arrays (FPGA), pp. 116--126 (2021). \doi{10.1145/3431920.3439301}

\bibitem{bayraktar2023cuquantum}
Bayraktar, H., Charara, A., Clark, D., Cohen, S., Costa, T., \emph{et al.}: cuQuantum SDK: A high-performance library for accelerating quantum science. In: Proc. IEEE International Conference on Quantum Computing and Engineering (QCE), pp. 1050--1061 (2023). \doi{10.1109/QCE57702.2023.00119}

\bibitem{xilinx2023alveo}
AMD: Alveo U55C Data Center Accelerator Cards Data Sheet, DS978 (v1.3) (2023). \url{https://docs.amd.com/r/en-US/ds978-u55c}

\bibitem{vitis2023hls}
AMD: Vitis High-Level Synthesis User Guide, UG1399 (2023.2) (2023). \url{https://docs.amd.com/r/2023.2-English/ug1399-vitis-hls}

\bibitem{xrt2023}
AMD/Xilinx: Xilinx Runtime (XRT) documentation (2023). \url{https://xilinx.github.io/XRT/}

\bibitem{nvidia2024cudaq}
NVIDIA Corporation: CUDA-Q: an open quantum--classical programming platform (2024). \url{https://nvidia.github.io/cuda-quantum/}

\bibitem{qiskit2024}
Javadi-Abhari, A., Treinish, M., Krsulich, K., Wood, C.J., Lishman, J., Gacon, J., \emph{et al.}: Quantum computing with Qiskit. arXiv:2405.08810 (2024)

\bibitem{nielsen2010quantum}
Nielsen, M.A., Chuang, I.L.: Quantum Computation and Quantum Information, 10th Anniversary Edition. Cambridge University Press (2010). \doi{10.1017/CBO9780511976667}

\bibitem{grover1996fast}
Grover, L.K.: A fast quantum mechanical algorithm for database search. In: Proc. 28th Annual ACM Symposium on Theory of Computing (STOC), pp. 212--219 (1996). \doi{10.1145/237814.237866}

\bibitem{boixo2018characterizing}
Boixo, S., Isakov, S.V., Smelyanskiy, V.N., Babbush, R., Ding, N., Jiang, Z., Bremner, M.J., Martinis, J.M., Neven, H.: Characterizing quantum supremacy in near-term devices. Nature Physics \textbf{14}(6), 595--600 (2018). \doi{10.1038/s41567-018-0124-x}

\bibitem{peruzzo2014variational}
Peruzzo, A., McClean, J., Shadbolt, P., Yung, M.-H., Zhou, X.-Q., Love, P.J., Aspuru-Guzik, A., O'Brien, J.L.: A variational eigenvalue solver on a photonic quantum processor. Nature Communications \textbf{5}, 4213 (2014). \doi{10.1038/ncomms5213}

\bibitem{markov2008simulating}
Markov, I.L., Shi, Y.: Simulating quantum computation by contracting tensor networks. SIAM Journal on Computing \textbf{38}(3), 963--981 (2008). \doi{10.1137/050644756}

\bibitem{vidal2003efficient}
Vidal, G.: Efficient classical simulation of slightly entangled quantum computations. Physical Review Letters \textbf{91}(14), 147902 (2003). \doi{10.1103/PhysRevLett.91.147902}

\bibitem{gottesman1998heisenberg}
Gottesman, D.: The Heisenberg representation of quantum computers. arXiv:quant-ph/9807006 (1998)

\bibitem{aaronson2004improved}
Aaronson, S., Gottesman, D.: Improved simulation of stabilizer circuits. Physical Review A \textbf{70}(5), 052328 (2004). \doi{10.1103/PhysRevA.70.052328}

\bibitem{chen2018quantum}
Chen, Z.-Y., Zhou, Q., Xue, C., Yang, X., Guo, G.-C., Guo, G.-P.: 64-qubit quantum circuit simulation. Science Bulletin \textbf{63}(15), 964--971 (2018). \doi{10.1016/j.scib.2018.06.007}

\bibitem{cross2017openqasm}
Cross, A.W., Bishop, L.S., Smolin, J.A., Gambetta, J.M.: Open quantum assembly language. arXiv:1707.03429 (2017)

\bibitem{ali2026multifpga}
Ali, N.: Multi-FPGA quantum circuit simulation: a distributed statevector architecture on four Xilinx Alveo U55C accelerators. Under review, The Journal of Supercomputing. Research Square preprint (2026). \doi{10.21203/rs.3.rs-10128153/v1}

\end{thebibliography}
